% =====================================================================
%  Harald Höffding, PERSONLIGHETSPRINCIPEN I FILOSOFIN.
%  Föreläsningar vid Helsingfors universitet våren 1911.
%  Auktoriserad öfversättning af Irene Leopold.
%  Stockholm: Wahlström & Widstrand, 1911.  Helsingfors: Aktieb. F. Tilgmanns
%  Bok- & Stentryckeri, 1911.  70 pp.
%  DIPLOMATIC TRANSCRIPTION OF THE SWEDISH PRINT.
%
%  THE TEXT.  Höffding gave these six lectures in Danish at Helsingfors in the
%  spring of 1911; his Danish manuscript is not known to survive, and the lectures
%  were published only in this Swedish translation, which he authorized.  (The
%  Danish of the 2013 Scientia Danica edition is Mogens Blegvad's retranslation
%  from this Swedish.)  The Swedish is therefore the earliest and the authoritative
%  text, and the English in translation.tex is made from it.
%
%  WITNESSES.  The Google Books scan of the University of Minnesota copy
%  (~/bibliotek/Høffding, Harald/Personlighetsprincipen_i_filosofin.pdf, 600 dpi).
%  (1) A tesseract (swe) reading of it, from which the draft was made
%  (ebook-method/personlighedsprincippet/ocr.py, draft.py); (2) Google's own OCR
%  layer in the same PDF, an independent reading, against which the draft was
%  collated word by word with punctuation (ebook-method/collate/check.py); every
%  disagreement was decided at the image (decisions.tsv there).  Print governs;
%  misprints are kept and marked PRINTED AS IS.  Spelling is that of 1911
%  (hvarje, lif, af, öfver ...) and is not modernized.
%
%  DONE 2026-10-03.  Draft (tesseract) collated against Google's layer: 218
%  disagreements, all decided: 73 at the image, 123 by class (the quotation marks
%  and dashes Google's layer drops, note anchors, plain tesseract slips), 22 by
%  check.py's rules; 66 of 66
%  page markers exact (check.py --pages).  Read by hand afterwards: the 13 footnotes
%  whole at the image (their italics -- the print has italics only in the notes;
%  the body has neither italics nor letterspacing, by a slant and a gap measure),
%  the note labels, and a reading of the whole text, which found some twenty
%  slips of the draft that the collation had passed (stray marks and dashes,
%  words joined or split; the rule that waves through dash-only differences let
%  "lik- som" through).  Misprints kept: làboratoriets
%  (p. 14), geründet (p. 16), the turned opening marks «la (p. 25) and «historien«
%  (p. 47), the missing opening mark before Justice (p. 49), enskida (p. 40),
%  erfarenhet (p. 43), and möjliggör (p. 19) where the sense needs omöjliggör.  Quotation marks as printed: »...» in the body, »...« in the notes.
%
%  CONVENTIONS.  "% --- p. N ---" + \opage{N} at the first word of printed page N
%  (a word broken over the turn is split with %).  Footnotes as printed, numbered
%  per page (\footnote[k]); a note that runs over to the next page's foot is given
%  whole at its anchor.  Italics \textit; letterspacing \emph.  Folios, the
%  footnote rule and the scanner's watermark are not transcribed.
% =====================================================================
\documentclass[11pt]{article}
\usepackage[utf8]{inputenc}
\usepackage[T1]{fontenc}
\usepackage{libertinus}
\usepackage{libertinust1math}
\usepackage[swedish]{babel}
\usepackage[protrusion=true,expansion=true]{microtype}
\usepackage{setspace}
\usepackage[margin=1.3in]{geometry}
\usepackage{fancyhdr}
\usepackage{hyperref}
\hypersetup{pdftitle={Personlighetsprincipen i filosofin}, pdfauthor={Harald Höffding}, hidelinks}
\pagestyle{fancy}
\fancyhf{}
\fancyhead[LE,RO]{\thepage}
\fancyhead[C]{\textit{Harald Höffding: Personlighetsprincipen i filosofin (1911)}}
\setstretch{1.3}
\usepackage{marginnote}
\newcommand{\opage}[1]{\marginnote{\footnotesize\textit{[#1]}}}

\begin{document}

% --- front matter: title (PDF 7) ---
\begin{center}
{\LARGE PERSONLIGHETS-}\\[0.3em]
{\LARGE PRINCIPEN I FILOSOFIN}\\[1em]
\rule{1.5cm}{0.4pt}\\[1em]
FÖRELÄSNINGAR\\
{\small VID}\\
HELSINGFORS UNIVERSITET VÅREN 1911\\
{\small AF}\\
{\large HARALD HÖFFDING}\\[1em]
\rule{1.5cm}{0.4pt}\\[1em]
AUKTORISERAD ÖFVERSÄTTNING\\
{\small AF}\\
IRENE LEOPOLD\\[6em]
\rule{3cm}{0.4pt}\\[4em]
STOCKHOLM 1911\\
WAHLSTRÖM \& WIDSTRAND
\end{center}
% (verso, PDF 8: HELSINGFORS 1911 / AKTIEB. F. TILGMANNS BOK- & STENTRYCKERI)
\clearpage

% ===== TEXT (pp. 5--70) =====
%
% --- p. 5 ---
\opage{5}
\begin{center}\large I.\end{center}

Med personlighetsprincipen förstår jag tanken, att människan på hvarje punkt och i hvarje afseende behandlas --- af andra och af sig själf --- som mål, aldrig blott som medel. Äfven där människan tjänar, --- tjänar andra människor, tjänar samhället, tjänar en sak, en idé, --- skall hon själf tillika vara mål, i det tjänandet bör leda till utveckling och förädling af hennes eget väsen, af dess individuella egendomlighet. Det är ett ideal som uppställes i denna princip, hvilken till först med full klarhet och bestämdhet formulerades af Immanuel Kant (1785). I de nordiska länderna gjordes den med stort eftertryck gällande af Erik Gustaf Geijer i Sverige (1842), från hvilken benämningen »personlighetsprincip» härleder sig. Medan Kant uppställde denna princip i motsats dels till auktoritets%
% --- p. 6 ---
\opage{6}
principen, dels till en yttre nyttighets- och lyckomoral, ställde Geijer den emot den romantisk-konservativa »historiska skolan», som enbart lade vikt på institutioner och traditioner, sådana dessa nu engång hade utvecklat sig, och han fann mer och mer att i denna princip innefattades allt hvad han litterärt, politiskt och religiöst hade verkat för. I Danmark uttalas principen, särskildt på det religiösa området, ungefär samtidigt (1846) af F. C. Sibbern och Sören Kierkegaard. --- Om vi följde principens historia längre tillbaka, skulle vi å ena sidan komma till kristendomen och den vikt den lade på den enskilda själen, å den andra till det sokratiska inskärpandet af själfkännedomens och själfuppskattningens betydelse.

Argumenteringen för principen är mycket olika hos dess skilda representanter. Till detta återkommer jag senare, då jag tillika skall undersöka huruvida denna princip öfverhufvud låter sig bindande bevisas. Här vill jag försöka framlägga allt hvad den innehåller. Den är först och främst uppställd som en etisk och religiös princip. Men jag tror, att den också innehåller synpunkter af stor bety%
% --- p. 7 ---
\opage{7}
delse för en rätt uppfattning af vetenskapens förhållande till lifvet. Den har intellektuell, icke blott etisk och religiös betydelse, och jag skall undersöka den från alla dessa tre sidor. Hvad Geijer sade om sin tid, att den »är en tid som lider af sina egna tankar», i det »den inre verlden svigtar i sina grundfästen, i sin tro, sin öfvertygelse, sitt vetande», --- och att »den grund- och modertanke, som här vill ut och gör sig väg, och under hvars födslovånda verlden bäfvar», just är personlighetsprincipen, %
% footnote 1
\footnote[1]{Samlade Skrifter (1853) I 6. p. 287.} --- det gäller säkert ännu den dag som är. Krisen är icke öfverstånden; den har blott antagit något andra former än på Geijers tid. I Geijers ord är träffande utsagdt, att här är en tanke som ofrivilligt arbetat sig fram, men som först genom att bli föremål för klart medvetande kan komma till sin fulla rätt.

Då vi nu begynna med personlighetsprincipens intellektuella betydelse, möter oss invändningen, att det i vetenskapen %
% --- p. 8 ---
\opage{8}
just gäller att se bort från vår personlighet, dess sträfvan, dess lust och smärta, dess hopp och dess fruktan. Böra vi icke just afkläda oss vår personlighet, om vi vilja vara forskare, om vi vilja finna tillvarons stora lagar, som gälla för allt och alla? I vetenskapen föras vi in i en samfäld värld. Men hvarje personlighet har sin egen lilla värld, olika för enhvar, liksom drömvärlden i motsats till den för alla gemensamma vakna verkligheten. Kan personligheten här fordra mera än att få lefva i sina drömmars värld, så länge vetenskapen icke drager också denna in under sin forsknings stränga rätt? ---

Jag går för ett ögonblick in på denna analogi mellan dröm och verklighet å den ena sidan, personlighet och vetenskap å den andra. Psykologin visar oss att samma lagar gälla för medvetandet under drömmen och i vaket tillstånd. Under båda uppenbarar sig ett behof och en förmåga att sammanfatta och ordna alla upplefvelser, alla förnimmelser, känslor och tankar. Det sker i drömmen på ett mera språngartadt och skenbart oregelbundet sätt, --- än genialt, än dumt. I vaket tillstånd sker det mera enligt %
% --- p. 9 ---
\opage{9}
bestämda teoretiska och praktiska synpunkter (ehuru också här både dumhet och genialitet kunna förekomma). Analogt är förhållandet mellan personlighet och vetenskap. Det är ett besläktadt behof och en besläktad förmåga som äro verksamma i båda. Om personlighet tala vi endast där en sträfvan efter sammanhang i själslifvet framträder, en sträfvan att samla alla psykiska elementer och leda dem i bestämda riktningar. Kontinuitet och koncentration äro personlighetens utmärkande kännetecken. Vetenskapen är blott en särskild, speciellt gestaltad form för denna egendomlighet. Försåvidt en motsats finnes mellan personlighet och vetenskap, bör det fastslås, att den eger rum inom personligheten, --- att den gäller förhållandet mellan personligheten som helhet och en af dess speciella lifsyttringar.

En af filosofins hufvuduppgifter är just att utforska, huru och hvarför det vetenskapliga medvetandet utvecklar sig som en speciell form för personligt lif. Det filosofiska eftertänkandet söker besvara frågan, huru den mänskliga andens natur yttrar sig, både i det arbete, som leder till en stor, gemensam världsbild, %
% --- p. 10 ---
\opage{10}
och i det arbete, som kanske blott får betydelse genom att bringa harmoni och ordning i den enskildes egen lilla värld. Denna lilla värld behöfver icke vara en drömvärld; här kan den ofvannämnda analogin icke utan vidare tillämpas; i hvarje fall finnes här ett särskildt problem, till hvilket vi senare skola återkomma.

Det är den kritiska filosofin som med största energi förfäktat, att tankens arbete i vetenskapen icke kan afslutas och icke kan förstås till sin fulla betydelse utan hänsyn till de medvetna väsendenas natur och villkor. Hvarje objektiv sanning har en subjektiv grundval, äfven om man icke kan sluta sig till denna sanning ur den subjektiva grundvalen i och för sig.

Vetenskapens historiska ursprung står att finna där, hvarest människan gjort den erfarenheten, att hon kan nå sina syftemål blott på alldeles bestämda vägar. Helst vill hon nå dem omedelbart. Endast tvungen af erfarenheten går hon omvägar. Och hon står då gentemot ett sammanhang, för hvilket hon måste böja sig. Här inträffar hennes första möte med den nödvändighet vetenskapen afslöjar på den ena punkten efter den %
% --- p. 11 ---
\opage{11}
andra. Vetenskapen fordrar resignation. Människan tänker här för att lefva. Tankens, och därmed personlighetens, arbete är här endast medel. Men det är ett medel som själft kan blifva mål, genom att bli föremål för omedelbart intresse. Det blir en själfständig uppgift att lära känna det stora sammanhanget, inom hvilket det personliga lifvet lefves, att finna en tingens ordning, som är mera omfattande än den enskildes egen lilla värld. Och då kan människan lefva för att tänka. Hon växer med de större syftemål hon sålunda uppställer för sig. Personligheten blir rikare, ju större krets af erfarenheter den omfattar, såframt den blott kan samla dem till harmonisk enhet, eller åtminstone ana och söka en dylik enhet. Det är lätt nog att koncentrera sig då det blott finnes ett fattigt och begränsadt innehåll att sammanfatta, då inga stora motsatser, ingen rikedom på upplefvelser lägga beslag på den psykiska energin. Personlighetens värld måste vidgas för att denna energi rätt må kunna träda i dagen.

Vetenskapen och dess historia bli först fattbara, då det finnes ett inre sammanhang mellan det vetenskapliga arbe%
% --- p. 12 ---
\opage{12}
tet och personlighetens hela lif och verksamhet. Det vetenskapliga arbetet förutsätter en energi och en hänförelse, en amor intellectualis, som icke är möjlig utan ett inre sammanhang mellan vetenskap och personlighet. Och ett sådant sammanhang förutsätter åter, att det icke är någon främmande makt personligheten böjer sig för då den genom forskningen ledes till att erkänna ett stort sammanhang, som visar långt ut öfver dess egen lilla värld. Längtan att finna sanningen på den väg vetenskapen anvisar, är blott en speciell form för längtan till öfverensstämmelse med sig själf under alla mångfaldiga och växlande upplefvelser och erfarenheter, --- den längtan, af hvilken personlighetens eget väsen består. Och till och med där personligheten skarpt förnimmer motsatsen mellan sin egen lilla värld och den oöfverskådliga, öfver alla mänskliga verk och syftemål nående värld vetenskapen öppnar blicken för, hör det just till personlighetens adel, såväl som till dess plikt, att med egen kraft, med ansträngande af sin egen tankes energi finna de stora lagar som betinga dess existens och dess värden. Personligheten luttras genom att se sin %
% --- p. 13 ---
\opage{13}
existens betingad och begränsad af tingens stora sammanhang.

Icke alla sidor af personligheten tagas i samma mån i anspråk af det vetenskapliga arbetet. I sin forskning måste människan utgå från bestämda synpunkter, följa bestämda vägar, utveckla sitt tankelif enligt vissa bestämda metoder, som visserligen äro speciella gestaltningar af personligt lif öfverhufvud, men som under loppet af mänsklighetens långvariga tankearbete bevisat sin oumbärlighet. Vetenskapen rör sig med vissa grundbegrepp eller grundpredikat, hvilka af gammalt kallas kategorier, och hvilka tillsamman utgöra en rustning, som lämpar sig såväl för personligheten som för det motstånd tankearbetet skall öfvervinna.%
% footnote 1
\footnote[1]{Jmfr. härom \textit{Den menneskelige Tanke, dens Former og dens Opgaver} (1910). Kap. III.} Man har stundom kallat personligheten, försåvidt den uppträder i denna rustning »det rena jaget», i motsats till »det empiriska jaget», som visar oss personligheten i hela dess individuella verklighet, med alla erfarenheter, föreställningar, önskningar och drifter. Det rena jaget är tankeensamhetens, studie%
% --- p. 14 ---
\opage{14}
kammarens och làboratoriets %
% PRINTED AS IS: „làboratoriets“ (p. 14)
jag, det med »kategorierna» strängt arbetande jaget. Men det är blott en särskild, genom arbetsfördelningen nödvändiggjord gestaltning af detta samma jag, som manifesterar sig i föreställningar, önskningar och drifter. Det är ingen absolut skilnad mellan dem, och deras inbördes förhållande kan framträda i otaliga nyanser hos olika individer. Det gifves forskare som förvandlas till helt andra människor, då de ur tankeensamheten träda ut i det vanliga lifvet. Men det är en ofullkomlighet som tyder på en olöst disharmoni i själfva det personliga lifvet.

På denna punkt står striden i våra dagar mellan kriticism och pragmatism. Den förra är (redan hos grundläggaren, Immanuel Kant) böjd att anse »det rena jaget», inbegreppet af den forskande tankens former (kategorierna) såsom gifna en gång för alla, oberoende af all erfarenhet och historia (också vetenskapens egen historia). Den senare är (både i den form Ernst Mach och den William James gifvit den) böjd att betrakta de former och förutsättningar, som äro nödvändiga för det vetenskapliga arbetet, såsom blott betingade af yttre uppgifter och prak%
% --- p. 15 ---
\opage{15}
tiska intressen, utan att de stå i något inre sammanhang med den mänskliga naturen själf. Enligt min uppfattning bortfaller denna motsägelse då man närmare undersöker de enskilda tankeformernas natur och utveckling. I min bok om »Den menneskelige Tanke» har jag arbetat i denna riktning, och har därvid blifvit styrkt i öfvertygelsen att människan i vetenskapen hvarken står blott som medel för vissa rena formers användande eller för tillfredsställandet af utifrån gifna uppgifter.

Då vetenskapen sålunda i sitt ursprung visar tillbaka till den mänskliga personligheten och framträder som en särskild form af det arbete hvarje personlighet har att utföra för att häfda och utveckla sig själf, blir dess personliga karaktär synlig från en annan och mera individuell sida vid afslutningen af det vetenskapliga arbetet. En absolut afslutning, en afslutning en gång för alla är icke möjlig, emedan lifvet och erfarenheten städse gå vidare. Men hvarje enskild forskare skall och måste stanna vid en bestämd punkt, där inre eller yttre orsaker hindra honom från att komma längre. Och då skall han själffallet söka en öfverblick, en samman%
% --- p. 16 ---
\opage{16}
fattning, en afrundning. Han skall tillika försöka en inblick i de outforskade, om också icke outforskliga regionerna, liksom man kastar ljus i mörkret med en elektrisk strålkastare. Här är det icke blott »det rena jaget» som arbetar; personlighetens öfriga sidor skola mer eller mindre medverka. Man ser det t. ex. på naturforskares försök att ställa sina resultat i förhållande till hela sin världsåskådning. Och man ser det när den historiske forskaren från kritisk granskning af källorna öfvergår till historieskrifning; denna kräfver en afrundning, vid hvilken forskarens personlighet gör sig mera gällande än under det kritiska arbetet. Man ser det ännu mera i filosofin, då den öfvergår från kritik till system, från analys till syntes.

Filosofiska system borde aldrig ha velat vara mera än ett individuellt afrundade. Afrundning kan vara berättigad, äfven om afslutning är omöjlig. Goethe karakteriserade sin forskning med principen »Nie geschlossen, oft geründet!» %
% PRINTED AS IS: „geründet“ (Goethe: gerundet) (p. 16)
 Och vid afrundningen bör den individuella karaktären uttryckligt erkännas. De verkligt betydelsefulla systemen förlora icke därpå, utan få just därigenom ett särskildt intresse. Till och med det system, %
% --- p. 17 ---
\opage{17}
som ofta anses såsom det mest objektiva och opersonliga af alla --- Spinozas --- visar sig vid närmare granskning genomträngdt af sin grundläggares personlighet; den utgör bandet som förenar de många tankeserierna till en enhet. Genom hela tankegången förnimmes mästarens oförskräckta konsekvens, hans stora resignation och hans innerliga fördjupande i tankearbetet, som för honom var en religion. Det är icke minst detta som gör att mången åter och åter, och med städse nytt utbyte vänder tillbaka till detta tankens konstverk. 

Det intresse filosofins historia erbjuder, beror till en stor del på detta personliga element. Filosofins historia kan betraktas som en stor dialog, hvars särskilda repliker härröra från forskarpersonligheter, hvilka hvar för sig, på grundvalen af sin tids vetenskap och kultur, ha utvecklat en tanketyp, en mer eller mindre väsentlig syn på tillvaron. För filosofin har dess historia, just på grund af personlighetsprincipen, större betydelse än för andra vetenskaper. Den står i detta, som i andra afseenden, emellan vetenskapen och konsten. Och för mången ligger just häri dess stora dragningskraft. %
% --- p. 18 ---
\opage{18}
\begin{center}\large II.\end{center}

För en rätt uppfattning af vetenskapen fordras att man i den ser en frukt af personligt arbete. Därigenom uteslutas såväl fruktan för dess resultat som förhäfvelse öfver dem.

Men personligheten är icke blott vetenskapens subjekt. Den är tillika själf ett af vetenskapens objekt. Vetenskapen vänder sig till den som till ett gifvet ämne, hvilket skall förstås och förklaras till sin natur och till sitt ursprung. Den personliga faktor, som medverkar i all vetenskap, kan och bör själf göras till föremål för vetenskap. Världen förstås igenom människan (ex analogia hominis); kunskapen om världen betingas af den mänskliga tankens förmögenheter och former. Men människan är själf en del af världen, som blott kan förstås i sitt sammanhang med den öfriga världen. Människan förstås alltså igenom världen (ex analogia universi).

Som Pascal sagt: »i rummet omfattar %
% --- p. 19 ---
\opage{19}
(comprend) världen mig, men i tanken omfattar (comprends) jag världen.»

De två betraktelsesätten strida icke mot hvarandra. Det är två synpunkter som båda måste anläggas, och som stå i inre växelverkan. Den mänskliga tanken, som arbetar på att lösa tillvarons gåtor, har själf utvecklat sig inom tillvaron. De förutsättningar, som betinga och begränsa vår världsuppfattning, ha själfva en historia. Det är icke minst nödvändigheten af denna dubbla synpunkt som möjliggör %
% PRINTED AS IS: „möjliggör“; the sense (and the next sentences) require „omöjliggör“ (p. 19)
en absolut afslutning af kunskapen. Vi kunna icke på en gång använda båda synpunkterna. Vi kunna icke på en gång stå på hufvudet och på benen; men vi kunna (kanske) göra först det ena och sedan det andra.

Mot att göra personligheten till objekt för vetenskapen har likväl gjorts en intressant invändning af en af det nutida Tysklands finaste tänkare, Heinrich Rickert, i hans skrift »Die Grenzen der naturwissenschaftlichen Begriffsbildung»(1901--- 1902). Rickert förfäktar en skarp motsats mellan naturvetenskap och historia (andevetenskap): den förra angår det universella och allmänna, det som ständigt upprepas; --- den senare har till före%
% --- p. 20 ---
\opage{20}
mål något som blott sker en gång (das Einmalige), enskilda händelser och individuella gestalter. Den enbart naturvetenskapliga metoden går ut på att finna allmänna begrepp och lagar, under hvilka de enskilda fenomenen kunna inordnas såsom exempel. De historiska vetenskaperna däremot stanna vid individuella totaliteter. De naturvetenskapliga begreppen äro allmänbegrepp, de historiska begreppen äro individualbegrepp!%
% footnote 1
\footnote[1]{»Von den Individuen gibt es keine Naturwissenschaft, weil sich nichts Besonderes mehr ihnen unterordnen lässt, im Vergleich zu dem sie noch als ein allgemeiner Begriff oder als eine »Natur« aufzufassen wären.» Über die Grenzen der naturwissenschaftlichen Begriffsbildung, p. 295.}

Mot denna uppfattning ställer jag påståendet, att i naturens, såväl som i andens värld äro allmänna begrepp och lagar slutligen blott medel att förstå det individuella och det historiska (das Einmalige). Det är vetenskapens stolthet, att den kan bilda begrepp och lagar, hvilka gälla för otaliga individer och händelser. Förmågan därtill var redan för Platon, den grekiska vetenskapens typiska gestalt, en gudomlig gåfva, som städse på nytt fyllde hans själ med beundran. Ja, %
% --- p. 21 ---
\opage{21}
enligt Platon voro själfva gudarna underordnade de eviga ideerna och gladde sig åt att se upp till dem. Men intet verkligt gifvet fenomen, hvarken i naturens eller i andens värld, blir uttömmande förklaradt genom att betraktas som exempel på ett allmänbegrepp eller på en allmän lag. För att förstå naturfenomenen såväl som de historiska fenomenen kräfves ett samspel af lagar. Lagarna äro de tankemedel, som genom sin samverkan göra ett fenomens uppkomst och utveckling fattliga. Solsystemet, jorden, organismen, personligheten, samhället --- det är exempel på individuella totaliteter, hvilkas historia det är vetenskapens uppgift att skrifva, med tillhjälp af de begrepp och lagar den förmår finna. Vetenskapens historia visar oss tydligt utvecklingsbegreppets stigande betydelse. Slutligen framstår uppgiften att skrifva »världens» biografi --- försåvidt »världen» kan te sig som ett afrundadt helt.

Det förhåller sig ingalunda så, att det blott är historien som uppvisar den enskilda företeelsen (das Einmalige). Icke häller i naturen förekommer absolut upprepande. Ingen tilldragelse återkommer någonsin på alldeles samma sätt. Natu%
% --- p. 22 ---
\opage{22}
rens oföränderlighet är blott skenbar. Det är ett stort världsdrama (eller kanske snarare ett epos), som naturvetenskap och historia i förening skola tyda. Däri förekommer ingen scen, ingen situation, ingen replik mer än en gång, ifall man noga lägger märke till alla enskildheter, alla nyanser och accenter.

Det är ett stort och upphöjdt ideal som härmed ställes för vetenskapen. Men det är nödvändigt för att rätt förstå vetenskapens betydelse. Naturvetenskapen närmar sig det medels de stora hypoteserna, som resas på grundvalen af funna lagar. Lagarna gälla städse enkla och elementära förhållanden; sedermera framställer sig uppgiften huru dessa förhållanden göra sig gällande inom de helheter af olika art och olika grad som erfarenheten uppvisar. Under det naturvetenskapen går från elementerna till totaliteterna, står andevetenskapen mera direkt inför totaliteter (tilldragelser, personligheter, samhällen), och det sådana, som icke låta upplösa sig i elementer lika lätt som de materiella totaliteter naturvetenskapen behandlar. Detta betingar dock blott en gradskilnad, ingen artskilnad. %
% --- p. 23 ---
\opage{23}
Det är af gammalt erkändt, att ju mera omfattande våra begrepp äro, desto innehållsfattigare äro de, och omvändt. Men det är icke innehållsrikedomen som skall underordnas de omfattande abstraktionerna; det är de senare som skola användas till att förstå och värdesätta den förra. De mest allmängiltiga vetenskaperna, logik och matematik, användas (oafsedt deras eget själfständiga intresse) såsom medel och förutsättningar för de mera speciella vetenskaperna. Svårigheterna ökas och gåtorna bli flere ju mera vi närma oss de individuella totaliteterna. Men dessa ange städse forskningens yttersta föremål och uppgifter.

Då det säges, att vi icke kunna bilda oss egentliga begrepp om det individuella (das Einmalige), emedan det icke gifves något som kan inordnas däri, måste därtill svaras, att till det begrepp vi bilda oss om ett individuellt fenomen, en tilldragelse (t. ex. den franska revolutionen) eller en historisk personlighet (t. ex. Luther) höra alla de olika sidor, egenskaper, situationer och utvecklingsgrader som kunna bidraga till dess fullständiga karakteristik. Också de »historiska» begreppen, icke blott de naturvetenskapliga, %
% --- p. 24 ---
\opage{24}
omfatta då en mångfald. Det måste göras åtskilnad mellan tvenne slag af individualbegrepp. Tänker jag mig en historisk personlighet i en bestämd situation (t. ex. Luther i Worms) kan jag om honom bilda mig ett konkret individualbegrepp; det får en åskådlig karaktär och bär minnesbildens eller fantasibildens prägel. Men ingen historisk personlighet uttömmes i en enstaka situation (lika litet som i en enstaka egenskap eller vid en enstaka utvecklingsfas). Mitt fullständiga begrepp om den måste motsvara alla de situationer i hvilka den uppträder, och det kan då kallas ett typiskt individualbegrepp. Den stora svårigheten med individualbegrepp är just att afgöra förhållandet mellan det konkreta och det typiska. Är det den unge Luther, han som brände påfvebullan, eller är det den gamle Luther, organisatören och dogmatikern, som skall bli afgörande för vårt slutliga begrepp om Luther? Hvilken af faserna i den franska revolutionen skall för oss bestämma dess (typiska) begrepp?

Lösningen ligger säkert däri, att det individuellas typiska begrepp anger den lag, till följd af hvilken öfvergången sker %
% --- p. 25 ---
\opage{25}
från det ena stadiet och den ena situationen till de andra. Det är förändringens lag som ger uttryck åt en individualitets väsen. Kan denna lag upptäckas, skall individens natur kunna återfinnas i hvarje af dess stadier, och man skall i det föregående stadiet kunna spåra de följande. %
% footnote 1
\footnote[1]{\textit{Leibniz}, individualitetsbegreppets egentliga grundläggare, har redan uttalat denna tanke. Det är «la %
% PRINTED AS IS: „«la“, the opening mark turned (p. 25)
loi de l'ordre«, la loi du changement«, som utgör hvarje väsens individualitet. Lettre à Basnage 1698. (Opera philos. ed. Erdmann, p. 151).} Hvarje psykologisk och historisk framställning arbetar på att finna en dylik lag. Hvarje god biografi låter oss i alla fall ana den. 

Konsten kan endast visa oss det konkreta, bilder af enskilda situationer och handlingar, i hvilka en personlighet träder i dagen, icke hela den följd af bilder som vore af nöden för att bilda ett typiskt begrepp. Därför kunna stridiga tolkningar uppstå om konstnärligt framställda gestalter, t. ex. om Hamlet. Och dock förmår konsten genomföra individualiseringen bättre än naturen och historien i regeln göra det.

Så vanskligt det kan vara att bilda typiska individualbegrepp, framstå dessa %
% --- p. 26 ---
\opage{26}
dock slutligen som tankens högsta uppgift. För att kunna finna allmänna lagar, måste vetenskapen ofta bortse från det faktum, att allt hvad erfarenheten visar oss har en individuell karaktär, framträder som en helhet, låt vara enkel och elementär. Det gifves många grader af individualitet --- ända från ett moln på himlen eller en sandhög på stranden till en mänsklig personlighet eller ett mänskligt samhälle. Alltför ofta träffar man den föreställningen, att naturen »börjar» med ett kaos af elementer och krafter utan kvalitet och bestämd riktning, och att alla totaliteter och individualiteter äro deras produkter. I våra läroböcker börja vi med det enkla och öfvergå till det sammansatta, emedan förståendet lättast uppnås på detta sätt. Men hvad vi veta om elementerna, om det enkla är i själfva verket alltid vunnet genom analys af totalsammanhanget, genom abstraktion från individuella fenomen, som kunna variera i oändlighet. %
% footnote 1
\footnote[1]{Jmfr, med afseende på den fråga som behandlats i detta afsnitt, min uppställning af »Totalitet och utveckling« som särskilda kategorier. »Den menneskelige Tanke«, p. 218--237.}

Den här anställda betraktelsen får en %
% --- p. 27 ---
\opage{27}
särskild betydelse, då personligheten blir objekt för vetenskapen. Ty personligheten är icke blott det oss närmast liggande exemplet på individuell totalitet. Den är tillika det mest utpräglade exemplet på en dylik, i det den uppvisar den innerligaste samverkan af elementer och krafter, som ligger inom vår erfarenhet.

Vi vända oss nu till psykologin, erfarenhetsvetenskapen om personligheten, för att se hvad den kan lära oss om personlighetens väsen i sig själft. %
% --- p. 28 ---
\opage{28}
\begin{center}\large III.\end{center}

I alla erfarenhetsvetenskaper komma de afgörande definitionerna i slutet, icke i början. Psykologin arbetar sig långsamt fram mot ett personlighetsbegrepp. Alla dess olika metoder --- iakttagelse, beskrifning, jämförelse, experiment --- leda direkt eller indirekt till detta mål. Då vi redan nu försöka uppställa --- och i det föregående hafva förutsatt --- ett personlighetsbegrepp, kan detta blott ske på hypotesens väg.

Såsom ett karakteristiskt drag för personliga väsen måste främst enhetlighet i alla tillstånd och företag framhållas. Vi kalla människan i desto mera utpräglad grad en personlighet, ju mera alla hennes tankar, känslor och sträfvanden äro riktade mot ett mål, --- ju mera koncentration det finnes i hennes andliga lif, --- ju större kontinuitet det finnes i hennes %
% --- p. 29 ---
\opage{29}
sträfvanden. Trots alla förändringar och växlande situationer måste hos henne framträda en gemensam lag, en enhetsprägel, eller åtminstone ett sammanhang. På det mest upphöjda sätt visar sig detta då människan är underkastad stora svingningar, stadd i starka stridigheter, och hon likväl icke i de mörka stunderna förnekar hvad som på lifvets höjder stod för henne som sannt och rätt.

Personligheten vidmakthålles --- liksom allt själslif --- blott genom ett oafbrutet andligt arbete, hvari den psykiska energin yppar sig. Det är en ständig kamp att sammanhålla och harmoniera en gifven mångfald af förnimmelser, tankar, känslor, drifter och intressen, som gå i olika riktningar eller stå likgiltiga gentemot hvarandra. Ju större denna mångfald och olikhet är, desto mera andligt arbete kräfves det för att själslifvet skall bestå. Och ju innerligare sammanhanget och koncentrationen äro, desto större psykisk energi måste på detta arbete användas.

Det var detta begrepp jag lade till grund, då jag för omkring 30 år sedan utarbetade min framställning af psykologin. Det hade utvecklats hos mig %
% --- p. 30 ---
\opage{30}
genom kritik af den äldre engelska psykologin (den s. k. »associationspsykologin») som sökte förklara själslifvet --- och därmed personligheten --- som produkt af en mekanisk hopfogning af psykiska elementer (själfständiga förnimmelser, föreställningar o. s. v.), hvilka stode i ungefär samma förhållande till personligheten som en stenhög till det hus som uppföres. Uppgiften blef nu att visa, att det aldrig, så länge själslifvet eger bestånd, helt fattas sammanhang och enhet däri, ehuru dessa kunna antaga mycket elementära former. Jag stödde mig här särskildt på läran om sinnessjukdomarna, som bevisar, att själslig sjukdom isynnerhet består i upplösning af sammanhanget inom själslifvet, i försvagande af den psykiska energin. Senare bekräftades min uppfattning af Pierre Janets verk, särskildt af hans beskrifning af »psykastenin», en andlig svaghet, som karakteriseras däraf att intet psykiskt innehåll, inga psykiska elementer bortfallit, men att det fattas energi till sammanhållning och koncentration, hvarför patienter, som lida af denna sjukdom, icke förmå träffa afgöranden, fatta beslut. Det fattas hvad Janet träffande kallar »psykisk spännkraft» %
% --- p. 31 ---
\opage{31}
(tension psychologique), som motsvarar det jag kallar syntes eller psykisk energi.

Hos barnet är den psykiska energin ännu icke så utvecklad, att en stark och varaktig koncentration kan bli möjlig. Andlig utveckling består icke blott i ett mångfaldigt innehåll, som upptages genom erfarenhet och tradition, men först och främst genom den energi som åstadkommer ett inre sammanhang. Drömtillståndet karakteriseras isynnerhet däraf, att det saknar fast sammanhang och konsekvens i medvetandets innehåll. Det har, liksom barnets medvetande, en sporadisk karaktär. Däri öfvergår lätt en föreställning till en annan; situationerna växla utan att det i regeln väcker förvåning och utan att en förklaring öfver denna växling sökes. I det vakna medvetandet fasthållas och följas bestämda riktningar, och förändringarna motiveras, eller, om motivering tyckes brista, bli de gåtor. Äfven om här blott är en gradskilnad, då det ju också gifves vakna drömmar, visar denna skilnad dock tydligt den psykiska energins betydelse, isynnerhet då vi hafva skäl att antaga att drömmedvetandet betecknar ett tillstånd mellan det medvetna och det omedvetna. %
% --- p. 32 ---
\opage{32}
Det är emellertid icke blott af svaghet eller bristande utveckling som  själslifvets särskilda elementer --- förnimmelser, tankar, känslor o. s. v. --- kunna uppträda själfständigt och utan närmare sammanhang. Vi kunna också genom uppmärksamhet och analys framhålla och isolera särskilda elementer och såvidt möjligt betrakta dem hvart för sig. Vi måste särskilja för att kunna uppfatta tydligt. Men dylika distinktioner äro mer eller mindre konstlade. De äro reflexioner, eftertankens verk, och få icke förvexlas med det den omedelbara upplefvelsen erbjuder. Ju mera vi förmå fördjupa oss i våra rent omedelbara och ofrivilliga upplefvelser, desto klarare visar sig vårt själslif förflyta som en kontinuerlig ström, hvari intet element kan isoleras från de andra utan att förändras till sin karaktär. Denna kontinuerliga beskaffenhet hos de omedelbara psykiska upplefvelserna har i nyaste tid på ett glänsande och eftertryckligt sätt ådagalagts af William James och Henri Bergson.

Slutligen visar också den experimentella psykologin åt samma håll. Dess egentliga grundläggare, Fechner, bevisade, %
% --- p. 33 ---
\opage{33}
att de särskilda sinnesförnimmelserna icke uppstå oberoende af hvarandra, utan äro inväfda i hvarandra, i det att efter en mycket stark påverkan af sinnesorganen fordras en starkare retning än vanligt för att en sinnesförnimmelse af en viss styrka skall uppstå. Också de såkallade kontrastverkningarna vittna om det stora sammanhanget till och med i vårt elementära själslif. Den gröna färgen framstår mera utpräglad (»mättad») efter den röda, som är dess komplementärfärg, än efter andra färger. Förnimmelsernas absoluta själfständighet är blott skenbar.

Och om vi betrakta saken från den fysiologiska sidan, ledas vi i samma riktning. Hurudant än förhållandet månde vara mellan själ och kropp, står det dock fast, att själslifvet utomordentligt intimt är fästadt vid nervsystemet. Men dettas betydelse ligger däri, att det är ett stort koncentrationsorgan, genom hvilket kroppens alla delar stå i nära förbindelse och växelverkan, så att därigenom ett slutet och koncentreradt uppträdande af organismen som helhet blir möjligt. Och ju högre och mera energiskt utveckladt själslif vi ha skäl att antaga hos ett lefvande väsen, desto större roll spelar %
% --- p. 34 ---
\opage{34}
nervsystemet och desto mera centraliseradt är det. Också i den yttre, sinliga företeelsen manifesterar sig själslifvets och därmed personlighetens enhetsprägel.

Det är blott hos de väsen, hvilkas själslif nått en högre grad af enhet och klarhet, vi tala om personlighet. Det personliga lifvet kan framträda i ofantligt många olika grader; till en stor del beror olikheten mellan de skilda individerna på en dylik gradskilnad i koncentrationens intensitet. Blott hos människor talar man om personlighet. Djurets själslif måste, enligt alla erfarenheter, antagas vara dels så enkelt, dels så växlande och osammanhängande, att det närmast kan föreställas vara analogt med vårt drömlif.

Vi ha i det föregående definierat personlighetsbegreppet blott från den formella sidan. Vi ha talat om energi och koncentration, om enhet och sammanhang. Men också innehållet, som samlas till enhet, har sin betydelse. Två personliga väsen med samma förmåga af koncentration kunna likväl vara mycket olika, på grund af det olika innehåll hvari de uppgå och kring hvilket de samla sig, samt den olika riktning i hvilken deras energi utlöses. %
% --- p. 35 ---
\opage{35}
I reellt afseende får personligheten sin prägel af att den omfattar ett centralt innehåll, ett herrskande intresse och sträfvande som betingar den riktning i hvilken koncentrationen sker och den psykiska energin användes. Hvarje särskild upplefvelse, tanke, stämning, beslut får inom personligheten sin bestämda plats genom sitt förhållande till ett grundvärde, som utgör den individuella personlighetens reella särtecken, liksom graden af koncentrationens energi utgör dess formella särtecken. Den lilla värld hvarje personlighet omfattar, har för hvarje tid, på hvarje utvecklingsgrad en medelpunkt, i förhållande till hvilken alla andra elementer i det personliga lifvet bestämmas till värde och betydelse. Att känna sig själf är att blifva medveten om detta centrala intresse, detta grundvärde, som utgör hvad man kunde kalla det reella jaget. Om det icke funnes dylika beständigt verksamma elementer inom  oss, skulle vi vara främmande för oss själfva. Det kan finnas beständiga elementer af rent periferiskt slag inom oss; men den djupare själfkännedomen förutsätter centrala elementer, och det gör därför ett komiskt intryck då Holbergs Jeppe igen%
% --- p. 36 ---
\opage{36}
känner sig själf på sin ihåliga tand. Ett dylikt periferiskt element kan icke bestämma alla andra elementers plats i den inre världsordningen. %
% footnote 1
\footnote[1]{Om olikheten mellan det formella och det reella jaget jmfr min »Psykologi« V B, 5.}

I det personliga lifvet uppstår en kris, om de centrala elementerna förändras. Det kan ske genom plötsliga katastrofer eller genom långa tiders tvifvel och oro; men det kan också ske långsamt och omärkligt, sålunda att vi först efteråt uppdaga, att det skett en väsentlig förändring med oss och inom oss. Det kan arbeta sig fram något inom oss, som alldeles ofrivilligt bestämmer vår sträfvan, men hvaraf vi först senare bli medvetna. Ett exempel, som här ligger nära till hands, är Geijers förhållande till den tanke han gaf namnet personlighetsprincip. Den uppstod hos honom vid en sen tidpunkt i hans lif, i hans 55:te år, men han upptäckte då, att han länge följt den utan att själf veta det. I ett bref som han skref på sin sista födelsedag (1847), yttrar han sig därom på följande sätt. %
% footnote 2
\footnote[2]{Samlade Skrifter I, 1. p. XXXVIII.} »Det ligger en grundtanke i min %
% --- p. 37 ---
\opage{37}
hela lefnad, som allt tydligare träder fram, och som ej är mitt verk; ty den har ledt mig och fört mig allt intill denna dag. Skulle jag gifva den ett namn, så är den ingen annan än den så mycket omtvistade personlighets-principen, hvars profet jag nästan mot min vilja blifvit. Jag menar dermed, att jag från mig afhållit och afskuddat allt, som kunnat hindra och hämma utvecklingen af mitt eget innersta väsende. Det har mest skett utan mitt eget medverkande. Men detta har nu kommit öfver mig, och det blifver mig klart, att hvad jag litterärt, politiskt och religiöst har fäktat för, egentligen är denna min innersta personlighet och dess rätt, i den bästa mening.»

Sålunda erbjuder själfva personlighetsprincipens utvecklingshistoria ett exempel på en hufvudpunkt i det personliga lifvets natur och villkor. Vi ledas till att se, att personligheten icke kan sägas vara på uttömmande sätt bestämd genom det, som vid en gifven tidpunkt träder i dagen i klart medvetande och med fullt förstående. Liksom man antager, att i det organiska lifvet kunna förekomma skenbart plötsliga förändringar (mutationer), genom hvilka nya typer kunna %
% --- p. 38 ---
\opage{38}
uppstå, på samma sätt kunna också på det andliga området genom en stilla tillvext förberedas ändringar i riktlinjerna, som för den ytlige betraktaren kunna synas oförklarliga, ehuru de ha sin orsak i det som försiggått i själslifvets innersta kärna. Vi ha här ett viktigt exempel på den olikhet mellan konkreta och typiska individualföreställningar som framställdes i föregående afsnitt. Det typiska begreppet af en personlighet kan ernås blott genom att följa den genom dess olika utvecklingsgrader, och det är ofta de senare stadierna som leda uppmärksamheten på de tidigare stadiernas innehåll och betydelse. Den beskrifvande psykologin har här ett stort fält för sig. Och det är att hoppas, att då sinnet för personligt lif, dess former och dess faser, utvecklas ännu mera än hittills, skall här kunna utföras ett betydelsefullt arbete. Hittills ha dogmatiska och doktrinära uppfattningar af olika slag hindrat ett oafbrutet arbete i denna riktning. Den beskrifvande psykologin närmar sig gränsen mellan vetenskap och konst; men genom klarhet i uppfattningen och genom användandet af jämförande metod kan den eröfra ny terräng för vetenskapen. %
% --- p. 39 ---
\opage{39}
\begin{center}\large IV.\end{center}

Från personlighetsprincipens intellektuella betydelse öfvergår jag nu till dess praktiska betydelse, som är den ursprungliga.

Hvarje princip uppställer en fordran på vår tanke eller på vår vilja. Personlighetsprincipen ställer fordringar i båda riktningarna. Jag har försökt visa hvilka fordringar den ställer på tankelifvet. Den praktiska fordran den ställer blir, såsom jag nu skall försöka visa, af stor betydelse för etik och religionsfilosofi.

Både hos Kant och Geijer, hos Kierkegaard och Sibbern uppställes personlighetsprincipen i kritisk och polemisk riktning emot samhällets och traditionens herravälde öfver de enskilda personligheterna. Den fordrar deras fria och själfständiga utveckling. Men den innehåller likväl i sig möjligheten att låta den %
% --- p. 40 ---
\opage{40}
sociala såväl som den individualistiska uppfattningen af det mänskliga lifvet och den mänskliga handlingen komma  till sin rätt.

Från en rent social synpunkt betonas, att individen alltid uppstår inom ett samhälle. En isolerad individ existerar icke. Själfva det individuella samvetet, det innersta hos människan, är en produkt af det sociala lifvet, ett socialt arf. De individuella dygderna äro socialt betingade. Och å andra sidan är individen med sin längtan och sin sträfvan försvinnande i förhållande till samhällets lif, som fortfar i otaliga släkten. Han är blott en droppe i hafvet. Hans lycka, hans kunskap, hans samvete få försvinnande ringa betydelse i motsats till släktets och samhällets stora och beständiga utvecklingsgång. Och den enda måttstocken på individens betydelse ligger i det bidrag han förmår lämna till denna utvecklingsgång.

Individualisten svarar å sin sida: hvar skulle samhällets betydelse sökas, om ej i den säkerhet och det stöd det ger de enskida %
% PRINTED AS IS: „enskida“ (p. 40)
personligheternas utveckling? Ett samhälles värde mätes efter graden, enligt hvilken det fyller denna uppgift. Detta %
% --- p. 41 ---
\opage{41}
är den enda möjliga måttstocken, redan af det goda skälet, att samhället består af individer och blott existerar i och genom dem.

Gör den sociala uppfattningen individen till medel för samhället, så gör den individualistiska uppfattningen samhället till medel för individen.

Denna motsägelse framträder icke så länge lifvet väsentligt lefves under instinkters och traditioners ofrivilliga ledning. Men hvarje kulturutveckling kommer till en punkt, där individen icke längre känner sig solidarisk med samhället, utan detta och dess kraf stå som blott tradition, blott auktoritet, kanske som rent yttre makt. Det är en sådan tid som, enligt Geijers ord, lider af sina egna tankar. Då ha --- just till följd af den sociala utvecklingen --- själfständiga personligheter utvecklat sig; men samhället vågar ännu icke fullt erkänna deras själfständighet, och de kunna icke erkänna samhällets värde. Det är egentligen först under sådana tider som filosofin uppstår. Filosofin får däraf i sitt ursprung lätt en individualistisk karaktär. Den grekiska filosofin, renässansens tänkande och det adertonde århundradets filosofi äro bety%
% --- p. 42 ---
\opage{42}
delsefulla exempel härpå. Det nittonde århundradets filosofi betecknar --- både i sin romantiska och sin positivistiska hufvudform --- en reaktion häremot, liksom också de stora arbetarassociationerna, hvilkas uppkomst hör till detta århundrades största händelser, ha en antiindividualistisk karaktär. Men problemet reser sig åter och åter, och vi ha gått in i det nya århundradet med det.

I detta problem måste en bestämd åtskilnad göras mellan det historiska och det filosofiska åskådningssättet. Det är ett faktum, att individen utvecklas inom samhället. Själfva hans längtan och hans sträfvan, hans etiska föreställningar och den kritik han utöfvar på traditionen kunna icke förstås, om man utgår blott från den enskilda individen, tänkt som isolerad. Individen är såtillvida en produkt af sociala orsaker. Och samhället sträcker sig långt ut öfver individen, både i framtid och förtid. Men följer däraf med nödvändighet en förpliktelse för individen att böja sig under samhällets lagar och verka för dess välfärd? Individen kan ju betrakta den sociala utvecklingen, hvars frukt han själf är, såsom ett blott faktum, ett resultat, som han, %
% --- p. 43 ---
\opage{43}
om han har andlig och fysisk makt därtill, nu använder i sina egna individuella syftemåls tjänst. Det tjänar till intet att tala om en social skuld, som den enskilde skall afbetala genom att verka för samhällets fortsatta utveckling; ty frågan gäller just hvarför individen skall erkänna denna skuld. Hvarför skulle individen icke anse sig själf som suverän?

Hvad individualismen i sina extrema former gör gällande, är egentligen det, att personligheten skall vara mål utan att vara medel. Mot personlighetsprincipen strider det dock icke att personligheten är medel, tjänar en sak, en idé, tjänar samhället --- om blott detta tjänande är medel till en högre personlig utveckling för individen själf. Människan vexer, såsom Schiller sagt, med sina större syften. Men dessa större syftemål kan hon blott finna i samhällslifvet. Likasom tanken växer blott genom att söka sanningen, det vill säga genom att finna så stort sammanhang som möjligt mellan så många erfarenheter som möjligt, så växer känslo- och viljelifvet blott genom de erfarenhet %
% PRINTED AS IS: „erfarenhet“ (p. 43; both readings agree)
som göras och det arbete som utföres genom deltagande i samhällets och släktets lif. Men vill%
% --- p. 44 ---
\opage{44}
koret härför är, att samhällets former och fordringar modifieras så, att det sociala arbetet --- afbetalningen på »den sociala skulden» --- kan tjäna till att krafter och möjligheter, som eljes skulle slumra oanvända hos individen, komma till utveckling. Personligheten kan göra sig till medel för samhället blott då samhället också gör sig till medel för personligheten. Möjligheten härtill är kärnan i den såkallade sociala frågan. I hela historien finna vi stora grupper af människor som väsentligt behandlades som medel, utan att samhället ansåg sig ha plikter gentemot dem. Alla berättigade sociala reformer gå ut på att upphäfva denna dualism af mål och medel. Samhället måste, om det skall fortfara att bestå, erkänna personlighetsprincipen, oaktadt denna ursprungligen blifvit uttalad från individualistiska ståndpunkter.

Icke alla former af individualismen förbise likväl denna åskådning. Max Stirner proklamerade (i sin bok »Der Einzige und sein Eigenthum», 1845) en lika abstrakt som absolut jagfilosofi, enligt hvilken hvarken inom eller utom oss funnes något, hvars kraf vi skulle erkänna. Så länge du tror på sanningen, %
% --- p. 45 ---
\opage{45}
säger han till och med, tror du icke på dig själf och är blott en tjänare. Hvarje rätt är en företrädesrätt. Jag erkänner ingen rätt, och jag bryr mig icke om att ha rätt, om jag blott har makt! --- Här göres samhället icke endast till medel; det försvinner egentligen helt och hållet. I motsats härtill stöder sig en individualism eller »anarkism» som Kropotkins på erfarenheten af den betydelse den ofrivilliga sammanslutningen har för alla lefvande väsen, och han har på ett intressant sätt skildrat den i sin bok »Mutual Aid» (1902). Kropotkin opponerar sig blott mot centralisationen, mot statens ingripande i de spontana samhällsdaningarna i mindre kretsar. I grannskapsförhållandet ordnar lifvet sig som af sig själft. Hvad Kropotkin icke har öga för är, att dessa ofrivilliga sammanslutningar --- såsom de forna gillena och skråen --- kunna leda till stillastående i förtryck och orättvisa, och att därför ett ingripande kan vara af nöden. Men han har rätt däri, att staten icke utan nödtvång bör ingripa i de mindre kretsar, i hvilka lifvet ovillkorligt rör sig --- af samma orsak som samhället icke bör ingripa i den enskilda personlighetens fria utveck%
% --- p. 46 ---
\opage{46}
ling, med mindre detta blir nödvändigt af hänsyn till andra personligheters fria utveckling. Det är i de enskilda individernas inre som de nya krafter och tankar ha sitt upphof, hvilka kunna föra samhällslifvet framåt. Framtidens grodd spirar där. Det är således för sin egen skuld samhället bör erkänna personlighetsprincipen. Dess verkliga makt visar sig i den kraft och själfständighet hvarmed de enskilda individerna kunna lefva sitt lif inom dess ram. 

Man träffar ofta hos moderna sociologer en viss mystisk uppfattning af begreppet samhälle, i det de antaga ett »socialt» ursprung för ideer och känslor. Om därmed blott menas, att ideer och känslor utvecklas hos de enskilda genom växelverkan med andra, är det alldeles riktigt. Men när »samhället» åberopas som orsak, är det mystik. Samhället som sådant kan lika litet frambringa rätt, moral och religion som det kan frambringa poesi och vetenskap. »Folkvisorna» äro icke diktade af »folket» som helhet, utan af enskilda sångare, hvilkas namn äro glömda, emedan deras verk strax blef allas egendom och emedan det, genom att gå genom mångas händer undergick för%
% --- p. 47 ---
\opage{47}
ändringar, som ofta gjorde det till något annat än det ursprungligen var. På tidigare kulturstadier glömmas de individuella initiativen lätt. Och det är just sådana kulturstadier sociologerna hälst sysselsätta sig med. På senare stadier spela de grundläggande personligheterna en större roll och glömmas icke så lätt. En appell till »samhället» som förklaring af ett fenomen är i själfva verket ett asylum ignorantiæ. %
% footnote 1
\footnote[1]{Det är i detta sammanhang af intresse att E. G. Geijers affall från den »historiska skolan«, som var den romantiska föregångaren af den moderna sociologin, just framkallades af hans historiska studier. (Saml. Skr. I, 8. p. 520). Det vill säga, han insåg, att det icke var någon förklaring af ett fenomen att åberopa sig på »samhället« eller «historien« %
% PRINTED AS IS: „«historien«“, the opening mark turned (p. 47)
en bloc. Och han insåg, »att historien ej kan ge, men väl modifiera principen för handling, och att denna tvärtom måste vara hvarje tidehvarfs egen insats.«}

Därmed skall det berättigade i att ställa samhälle och individ i motsats till hvarandra icke förnekas. Begreppet samhälle innebär, att medels växlande generationers lif och arbete hopsamlas ett kapital, en potentiell energi, som nedlägges i institutioner och traditioner, ja i själfva språket. Här är något som visar ut öfver den enskilde, --- tillbaka till %
% --- p. 48 ---
\opage{48}
förtiden och ut mot framtiden, --- något som icke tillhör någon som hälst enskild individ. Det är värden, som äro bundna, och som blott kunna frigöras genom enskilda individers arbete. En samhällsordning värdesättes icke blott enligt det kapital den sålunda förfogar öfver, utan också enligt det sätt, på hvilket den förmår göra det rörligt. Det kan finnas döda värden som få sken af lif genom sociala institutioner. Om värdena äro lefvande eller döda kan blott afgöras däraf, om de komma fria personligheters utveckling till godo eller icke. Ty de enskilda personligheterna äro såväl de centra, i hvilka lifvets värden röja sig, som de centra, från hvilka arbetet att finna och frambringa värden utgår.

Enligt hela denna åskådning stå samhälle och individ samtidigt som mål och medel för hvarandra.

Den högsta etiska tanke är den, att den enskilda personligheten skall utveckla sig så, att den just genom sin själfständiga och egenartade utveckling på bästa sätt, direkt eller indirekt, hjälper andra att utveckla sig själfständigt och egenartadt. Häri består den sanna rättfärdigheten, som Westermarck med rätta kallar »blom%
% --- p. 49 ---
\opage{49}
man af alla sociala känslor».%
% footnote 1
\footnote[1]{Justice, %
% PRINTED AS IS: no opening quotation mark before „Justice“ (p. 49)
 the flower of all moral feelings. The origin and development of moral ideas.« I p. 124. --- Jmfr min »Etik« III, 9; IX, 1; X, 4.} Den förenar harmoniskt själfhäfdelse och hängifvenhet, individens och samhällets intressen. Genom den lefver samhället i individerna, icke blott utom eller öfver dem. I den kommer personlighetsprincipen till sin fulla rätt. %
% --- p. 50 ---
\opage{50}
\begin{center}\large V.\end{center}

Det spörsmålet kan nu icke längre skjutas undan, om personlighetsprincipen, med hvilken vi här ha sysselsatt oss, kan bevisas vara giltig. I närmast föregående föreläsning har visserligen gifvits en motivering, men det kan bli fråga om icke denna motivering själf hvilar på vissa förutsättningar. Den bestod däri, att nämnda princip möjliggjorde den mest harmoniska utveckling af samhället och den innerligaste växelverkan mellan individerna. %
% footnote 1
\footnote[1]{Denna bevisföring gaf jag först i min Etik (VIII, 6 jmfr XIII, 6), där jag också kritiserat Kants konstlade argumentering. Geijers bevisföring (Saml. skr. I, 5, p. 132--137) stöder sig främst på, att människan i hvarje annan människa upptäcker ett väsen af samma värde som hon själf, samt på att människan blott kan utveckla sig till personlighet i växelverkan med andra människor. Men därmed uteslutes dock icke, att människan kan göra dessa besläktade väsen, hvilkas medverkan hon icke kan undvara, till blotta medel för sig själf! --- Karl Marx (Das Kapital I p. 646) och Leon Bourgeois (i La philosophie de la solidarité. Paris 1902, p. 38) uppställa satsen som själfklar, hvilket en blick på historien lätt visar vara oriktigt.} Men denna tankegång tages %
% --- p. 51 ---
\opage{51}
naturligtvis för giltig blott af dem, för hvilka en dylik utveckling af människosläktet står som ett godt. Den följer af tanken på universell mänsklig välfärd. Men frågan gäller huru det förhåller sig med själfva denna tanke. Svaret på denna fråga blir afgörande för etikens ställning som vetenskap.

Värde kan blott mätas med värde. Och värde visar tillbaka till känsla --- och om känslor, säges det, kan man icke disputera. Dock kan man, då ett visst värde kännes och erkännes, däraf sluta, hvilken plats andra värden konsekvent intaga i förhållande till det gifna värdet. Hos hvarje människa --- och i det närmaste hos hvarje folk och i hvarje tidsålder --- gör ett visst grundvärde sig gällande, ett intresse som ligger djupast och utöfvar det starkaste inflytandet. Det motsvarar, eller rättare tillhör hvad jag kallar personlighetens reella sida. Och det afgör hvad som för öfrigt får värde, och hvilket värde det får. Den, för hvilken egen fördel är grundvärdet, %
% --- p. 52 ---
\opage{52}
kan icke utan inkonsekvens göra uppoffringar till förmån för andra personer eller för allmänna syften. Omvändt kan den, hos hvilken människokärleken bestämmer grundvärdet, icke utan inkonsekvens handla egoistiskt. Våra etiska omdömen bestämmas af våra grundvärden, och dessa å sin sida af våra rådande känslor.

Men kan man nu påstå, att alla människor på alla tider och under alla förhållanden inneha samma grundvärden? Det vore helt visst ett opsykologiskt och ohistoriskt påstående, oaktadt tänkare som Kant och Bentham, hvar på sitt sätt, gjort det. Grundvärdena uppstå genom psykologiska och historiska processer, som aldrig föreligga helt afslutade. Ett grundvärde kan aldrig uppställas som själfklart. Det måste växa fram, och det kommer i hvarje enskildt fall an på huru långt växten framskridit. Det hjälper ej att säga, att människan växer med sina större syften, ty människan måste ha vuxit för att kunna uppställa för sig större syften och sedan åter växa vidare genom arbetet för dem. Det tjänar som Goethe sagt, icke till något, att tala till larven, som krälar i stoftet, om att den %
% --- p. 53 ---
\opage{53}
engång skall flyga från blomma  till blomma.

Personlighetsprincipen som allmän etisk princip förutsätter utvecklingen af ett humanitetsmedvetande, en allmän människokärlek, sådan den först framträdde under de sista århundradena före Kristus genom ett egendomligt sammansmältande af platoniska och stoiska ideer. Senare har den haft en lång och växlande historia. Den har haft att kämpa icke blott med individers egoism, utan också med familjers, klassers, nationers och konfessioners egoism. I det adertonde århundradet bröt den fullt igenom och ställde sina bestämda kraf på det borgerliga samhället. Men den är beständigt en stridande kyrka. Därför kunna de konsekvenser, som dragas ur det allmänna humanitetsmedvetandet, icke häller bli giltiga för alla. 

Vår tanke förmår icke annat än draga konsekvenserna af vissa bestämda förutsättningar. Det är erfarenhetens sak att afgöra om dessa förutsättningar finnas öfverallt och hos alla. Det förringar icke tankens, ej häller det etiska tänkandets betydelse. Tanken kan blott föra till klart medvetande det, som våra förut%
% --- p. 54 ---
\opage{54}
sättningar innebära. Och den filosofiska etikens betydelse beror också just af detta. Den skall framställa de konsekvenser, som med fog kunna dragas från ett visst stadium i det mänskliga känslolifvets utveckling. Senare kan undersökas huruvida konsekvenser från andra utvecklingsstadier sammanfalla med dem, och hvilken skilnad olika motivering och olika förutsättningar medföra!%
% footnote 1
\footnote[1]{Denna uppfattning af etikens vetenskapliga begränsning och betydelse har jag för längesedan utvecklat i min Etik (Kap. III). Det är mig en glädje att träffa på en liknande tankegång i Westermarcks »Origin and development of moral ideas« (I p. 17; 105; 122).}

Där den etiska motiveringen står vid sin gräns, emedan dess förutsättningar icke längre erkännas, uppstår frågan huru dessa förutsättningar kunna bringas till stånd. Detta blir moralpedagogikens sak. Det är genom uppfostran grundvärden bildas. Ordet uppfostran tages här i vidaste bemärkelse, sålunda att mänsklighetens hela historia ses som en stor uppfostran. Det är psykologins och sociologins sak att visa huru denna uppfostran ofrivilligt och omedvetet försiggår hos enskilda individer och i generationer, och %
% --- p. 55 ---
\opage{55}
hvilka medel som stå till förfogande då man fullt medvetet vill uppfostra människor.

Schopenhauer har sagt, att det är svårt att motivera moral, men lätt att predika moral. I den första delen af sin sats har han rätt. Men han har orätt i att det skulle vara lätt att predika moral, ifall därmed menas att föra människor till en sådan utveckling af deras personlighet, att de af egen kraft och åtrå kunna finna det rätta. Icke häller under uppfostran får ett personligt väsen betraktas och behandlas blott som medel, icke ens som medel för sin egen framtid. Larven bör behandlas som larv, ända tills den blir fjäril. Det är den stora tanken i Rousseaus pedagogik, att barnet skall behandlas som barn och icke blott som en blifvande vuxen. Det gäller således att taga den, som skall uppfostras, på den ståndpunkt där han står, och sedan på själfutvecklingens och själfverksamhetens väg leda honom till den nivå, det grundvärde, som kan bilda utgångspunkt för de slutledningar man vill att han skall draga. Här gäller Spinozas ord, att det icke gifves något sätt på hvilket man bättre kan visa sin andliga kraft och förmåga, än genom att föra andra %
% --- p. 56 ---
\opage{56}
till att lefva enligt deras egen klara öfvertygelse.

Äro således svårigheter förknippade med bevisledningen för etiska principer, finnas också svårigheter vid deras genomförande i hvarje enskildt fall. Äfven om grundvärdet förutsättes vara det samma, uppstår likväl frågan om man, utgående från de samfälda förutsättningarna, kan ställa samma etiska fordringar på alla individer. Individer äro olika, såväl i afseende å arten af sina förmögenheter som i afseende å den grad af psykisk energi de förfoga öfver. Däraf följer, att man icke kan fordra det samma af alla. Hvarken hvad som fordras eller huru mycket som fordras kan en gång för alla fastslås, utan måste bestämmas med afseende å de enskilda individernas natur och krafter. Samma börda tynger olika på olika skuldror. Det finnes människor som nödgas kämpa en hård strid för att uppfylla de mest elementära fordringar i afseende på själfbeherskning, dem andra fullgöra alldeles ofrivilligt. Det etiska arbetet kan vara det samma, ehuru resultaten blifva af högst olika värde. Människorna äro som löpare, hvilka begynna loppet från mycket olika utgångspunkter. %
% --- p. 57 ---
\opage{57}
Ty man kan icke neka till, att i etiskt afseende, liksom i andra afseenden, äro människorna mycket olika utrustade från början.

Det skulle strida mot personlighetsprincipen att dömma den enskilde efter en abstrakt regel, som icke gör afseende på hans individuella betingelser. Uppgiften blir att individualisera de etiska krafven, så att hvarje enskild kommer att utföra det arbete, som öfverensstämmer med hans personliga betingelser, och därigenom föres in på ett spår, som kan leda vidare. Från en sådan individualisering är etiken ännu långt aflägsen, och den fordrar ju också stor psykologisk urskiljning och stor praktisk människokännedom, såväl som stort tålamod. Det är lättare att systematisera och att uttala allmängiltiga satser, än att finna de individuella utgångspunkterna i de enskilda fallen. System och variation stå i etiken, liksom i biologin, i skarp motsats till hvarandra.

I begrepp, sådana som billighet och benådning är individualiseringen dock längesedan erkänd. Begreppet billighet uppställer redan Aristoteles som en säregen form för rättvisa, hvilken nödvän%
% --- p. 58 ---
\opage{58}
diggöres däraf, att det gifves individuella nyanser, som hindra tillämpandet af en allmän lag. Och att benådning tillämpas, kan rättsligt blott motiveras därmed, att lagstiftningen icke är i stånd att taga hänsyn till alla individuella förhållanden. Benådningen är, säger Jhering, rättvisans själfkorrigering.

Det är dock icke blott nedåt, som individualisering fordras, utan äfven uppåt. Det gifves icke blott individer, som börja sin etiska utveckling från en lägre nivå än genomsnittets; det gifves också individer som från begynnelsen äro utrustade med särskilda betingelser för etiskt arbete. Och i kraft af personlighetsprincipen måste den etiska lagen ställa större kraf på dem än på andra. De måste själfva fordra mera af sig själfva, än som fordras af andra. De förnedra sig själfva med att låta sig nöja med ett minimum. Det är deras »plikt» att inlägga »förtjänst». Schillers Wilhelm Tell erkänner detta, då han besvarar sin hustrus klagan öfver, att hans landsmän alltid vända sig till honom då det gäller något stort, med orden: »Ein Jeder wird besteuert nach Vermögen!»

I pedagogiskt afseende kan distinktionen mellan »plikt» och »förtjänst», som %
% --- p. 59 ---
\opage{59}
den katolska teologin lägger så stor vikt på, vara af betydelse. Dels kan det vara svårt nog att få alla att fullgöra de minimala fordringar man ofta tänker på vid ordet »plikt», dels kan det vara af nöden att varna för etiska vågstycken, för att våga sig längre än krafterna medge. Svårigheten ligger däri att vi blott genom försök, genom att våga lära känna våra krafters omfång. Genom att hålla oss tillbaka kunna vi förorsaka, att krafter och möjligheter hos oss, som kunnat utvecklas vid arbetet för större syften, bli obegagnade. Och vi lära icke känna oss själfva. Måhända leder försöket att utföra en stor handling till att vi lära känna vår begränsning. Måhända leder det rentaf till en tragisk utgång. Men sådana äro nu engång lifvets villkor. Som Longfellow sagt: »Blott de som trotsa hafvets faror, känna hafvets hemlighet.» %
% --- p. 60 ---
\opage{60}
\begin{center}\large VI.\end{center}

Det återstår ännu att undersöka personlighetsprincipens betydelse för  lifsåskådning och tro, alltså dess religiösa betydelse, då man tager ordet »religiös» i vidare bemärkelse. Jag tror att denna betydelse är i ständig tillväxt, och att den väsentligaste sidan af det religiösa problemets ställning i våra dagar beror därpå.

Det är tre elementer, som vid hvarje gifven tid bestämma människors tro: vetenskap, tradition och personlig erfarenhet. De spela en mycket olika roll hos olika individer, och till och med där de alla göra sig starkt gällande, förenas de inbördes på mycket olika sätt, hvarigenom lifsåskådningen får en skiljaktig klangfärg. Det finnes knappast två människor, hos hvilka de framträda på samma sätt eller förbindas på samma sätt. Det vore %
% --- p. 61 ---
\opage{61}
ett intressant studium att tränga in i dessa olikheter. Tyvärr är ett sådant studium blott i ofullkomlig grad möjligt gentemot samtida.

Mera och mera leda de två första elementerna, vetenskapen och traditionen, till det tredje, en personlig erfarenhet af lifvets gång och lifvets villkor.

a. Den nyare vetenskapen, hvars ursprung står att söka i renässansens århundrade, har isynnerhet genom två hufvudtankar utöfvat stort inflytande på lifsåskådningen. Den ena är föreställningen om universums oöfverskådlighet, om rummets och tidens obegränsade utsträckning. Den begränsade skådeplats, på hvilken människan dittills hade kunnat föreställa sig tillvarons epos utbredt, tog sig nu ut som en försvinnande punkt. Den andra är föreställningen om en stor lagbundenhet i hela naturen, en lagbundenhet, som steg för steg bevisas af att hvarje förändring i naturens eller andens värld förklaras genom sitt bestämda sammanhang med andra förändringar, utan att orsaker, som ligga utom en möjlig erfarenhet, påkallas. Detta vetenskapliga kausalbegrepp användes först på det oorganiska området, senare också på det %
% --- p. 62 ---
\opage{62}
organiska området, likaväl som på det andliga området.

De intryck denna stora förändring i världsuppfattningen gjorde på människors lifsåskådning voro mycket olika. Hos några väckte den entusiasm, i det tillvaron nu mera än förr fick det upphöjdas prägel, samtidigt som den öfvertygelsen klarnade, att det trots allt är samma krafter som verka i människans natur och i det oändliga universum. Hos andar som Bruno, Böhme och Spinoza gör denna riktning i lifsåskådningen sig gällande. Hos andra väckte denna utvidgning af horisonten, som förvandlade den jordiska skådeplatsen till en försvinnande punkt, fasa och yrsel, och ledde till ett lidelsefullare fasthållande af den gamla tron. Pascal är det stora exemplet på denna riktning. Det fanns slutligen dem, som alls icke lade vikt på den stora förändringen i världsbilden, icke blefvo medvetna om dess principiella betydelse, antingen de i sin lifsåskådning höllo sig till den traditionella tron eller gingo egna vägar. --- Ännu den dag som är finner man dessa tre ståndpunkter gentemot det stora problemet.

Man behöfver icke öfverdrifva veten%
% --- p. 63 ---
\opage{63}
skapens betydelse för att vara säker på, att den kommer att utöfva ett stigande inflytande på lifsåskådningarna. Man skall icke kunna gå ytterom vetenskapen. Redan Pascal sade ju, att om jorden verkligen roterar, så roterar den med både troende och icke troende. Mest skall vetenskapen verka genom den tysta kraft som ligger i dess metod, dess sätt att fråga och svara. Mindre och mindre skola människor, som uppfostrats med denna metod, kunna taga mystiska och antropomorfiska förklaringar som verkliga svar på erfarenhetens gåtor. Det är mindre vetenskapens speciella resultat, än det sätt att tänka den lär människorna, som skall bestämma lifsåskådningarnas framtid. I denna framställnings första afsnitt sökte jag dessutom framhålla, att det icke är någon främmande makt som här träder personligheten till möte. Det vetenskapliga intresset och den vetenskapliga förmågan hänga nära ihop med personlighetens väsen och bära dess prägel. Den djupa helhetslängtan som är det innersta hos personligheten, ligger också till grund för vetenskapen. Denna är personlighetens verk, som den icke kan förneka. Och det hjälper icke att sluta ögonen för den %
% --- p. 64 ---
\opage{64}
vidgade horisont vetenskapen öppnar. Vi måste se oss själfva och skådeplatsen för allt vårt arbete och alla våra intressen såsom mindre kretsar innanför den stora kretsen, hvars medelpunkt är öfverallt, och hvars omkrets är ingenstädes. Emellertid skall det knappast någonsin bli möjligt att, utgående från vetenskapens synpunkter och resultat, konstruera lifsåskådningar för de enskilda personligheterna hvar för sig. Såsom vi sett (i slutet af andra och tredje afsnittet), är det visserligen vetenskapens yttersta och högsta uppgift att förstå och förklara individuella totaliteter, alltså också personligheter, men här kan blott bli fråga om ungefärliga antagenden, och vetenskapen närmar sig här konsten. För den enskilde skola de rent individuella lifserfarenheterna städse bli det väsentliga vid den afrundning han ofrivilligt gör då lifsåskådningen bildas. En människa är mera än ett exempel på en allmän lag, och till och med det människolif, som rör sig i den trängsta krets, kan ge sitt bidrag till mänsklighetens stora samerfarenhet. I hvarje fall kan blott den enskilde veta hurudan lifvets och ödets gång är på den särskilda plats inom till%
% --- p. 65 ---
\opage{65}
varon, där hans existens befinner sig. »Det ringaste försvarar sin plats i lifvets krans», har en dansk diktare sagt.

b. På hvilket sätt de religiösa traditionerna än uppstått, bero deras fortsatta bestående och den betydelse de allt ännu äga för så många, på den omständigheten, att i dogm och kult nedlagts oräkneliga generationers djupaste  lifserfarenheter. Religionshistorien visar oss huru föreställningar, hvilka tidigare voro religiöst indifferenta, under inflytande af stora rörelser inom känslolifvet, blefvo tagna i religionens tjänst. Så har det t. ex. gått med idén om själavandringen i Indien och med den grekiska filosofins tankar i den kristna kyrkans dogmatik. Men de betraktas inom religionerna icke blott som symboler för lifserfarenheter. De befästas som dogmer, hvilka skola innehålla lösningen af lifvets gåtor. De ryckas upp ur den lefvande jordmån hvari de grott. Dock röjer sig mera och mera i vår tids religiösa medvetande en brytning mellan den symboliskt-poetiska och den bokstafligt-dogmatiska uppfattningen. Det framträder allt klarare, att de gåtor vetenskapen icke kan lösa, häller icke lösas genom de religiösa föreställ%
% --- p. 66 ---
\opage{66}
ningarna, redan af det skäl, att vetenskapen både frågar och svarar annorlunda än religionen. Den tid är förbi, då religionen kunde tillfredsställa alla sidor af det mänskliga själslifvet. Intellektuellt har den förlorat drabbningen, huru det sedan må förhålla sig med andra sidor af själslifvet. I våra dagar går mycken ädel kraft till spillo hos de mest själfständiga andar med det arbete brytningen mellan religion och vetenskap framkallar, --- kraft, som kunde användts till att grundlägga en på egen erfarenhet fotad lifsuppfattning.

Och dock är det icke förhållandet mellan vetenskap och religion som är det viktigaste vid det religiösa problemets ställning i våra dagar. Från båda sidor har sagts hvad som kan sägas, och fortsatt diskussion skall icke bringa något nytt. Den afgörande punkten ligger däremot i förhållandet mellan religion och personlighet. I motsats till den religiösa traditionens oföränderlighet står den fri, den personliga lifserfarenhet, hvars rikedom och fullhet stegras i den mån de enskilda personligheternas egenart och berättigande hvar för sig erkännas. Hvarje mänsklig individ är en af de medelpunkter, %
% --- p. 67 ---
\opage{67}
där lifvets värde skönjes, och en af de utgångspunkter, från hvilka nya värden kunna finnas eller frambringas. Det är icke så mycket vetenskapens som lifvets domstol, inför hvilken religionen i en framtid skall ställas. Det sker i kraft af personlighetsprincipen. Ty lika litet som en människa blott är ett exempel på ett vetenskapligt begrepp, lika litet skall hon vara blott medel till vidmakthållandet af en religiös tradition. %
% footnote 1
\footnote[1]{Då jag i det föregående ofta hänvisat till Erik Gustaf Geijer såsom en af personlighetsprincipens viktigaste representanter, skall jag anföra ett yttrande som visar, att han också i sin religiösa uppfattning blef sin princip trogen. »Jag är«, skrifver han i ett bref, »hvarken en kyrko-christen eller ens en bibel-christen . . . Jag är en christen på egen hand. Tusende befinna sig med mig i lika läge, i själfva verket alla som ej kunna nöja sig med en blott auktoritetstro och som äro nödgade att ur egen erfarenhet och eftertanka bygga sig sin egen religion. \textit{Detta individuella, nu det enda lefvande i religionen} . . . pekar dock öfverallt till ett gemensamt mål.« (Saml. Skr. I, 1. p. XXXII 7). (Kursiveringen af mig). Dock gaf Geijer, som stod under starkt inflytande af sin tids romantiska spekulation, tillika sin personlighetsprincip en teologisk och metafysisk betydelse, som är ohållbar då hvarje försök i denna riktning leder till ett spel med analogier och allegorier, som hvarken kunna konsekvent genomföras eller bekräftas af erfarenhet. Jmfr härom min Religionsfilosofi § 32.}

% --- p. 68 ---
\opage{68}
Det religiösa behofvet är ett behof att bevaka lifsvärdena ut öfver de gränser, inom hvilka den mänskliga viljan kan arbeta för dem. Kunde människan själf absolut säkerställa sin andliga och materiella egendom, skulle icke religion i någon form uppstå. Men nu uppkomma hopp eller fruktan, hänförelse eller nedslagenhet, allt efter det man känner sig på lifvets höjder eller i dess djup. De stora religionernas kulter gifva i den upphöjdaste poesi uttryck åt dessa motsatta stämningar, som lifvets epos för med sig och som alltid skola göra sig gällande hos enhvar, som lefver lifvet helt och fullt, och hvars sympati icke förslöas och begränsas. 

Jag vill icke spå om framtidens religion. Jag är filosof, icke profet. Men allt talar för, att ju mera de verkliga och beständiga lifserfarenheterna få inflytande på det religiösa lifvet, desto mera skall den symboliskt-poetiska uppfattningen af de religiösa föreställningarna aflösa den bokstafligt-dogmatiska, och göra det allt mera medvetet. Och därmed skall ett stort steg mot fri, själfförvärfvad mänsklig lifsåskådning tagas. ---

c. Både vetenskap och tradition visa %
% --- p. 69 ---
\opage{69}
således mot den personliga lifserfarenheten. En genomgående individualisering af lifsåskådningen är nödvändig. Personligheten skall sätta sin prägel på sin tro. Blott sålunda blir denna ett vittnesbörd om huru lifvet rör sig på denna särskilda plats i världen. Och då framskridna psykologiska undersökningar leda oss till att inse, att de individuella olikheterna äro långt större och djupare liggande än hittills antagits, skola sannolikt framtidens lifsåskådningar uppvisa större motsatser, än man kunde få blick för sålänge de samfälda konfessionella dogmerna täckte de inbördes afvikande uppfattningarna och användningarna af en lika lydande bekännelse.%
% footnote 1
\footnote[1]{Jmfr den öfversikt jag i min religionsfilosofi (§ 36--43 \& § 94--96) sökt gifva af individuella olikheter, som kunna bli af betydelse för lifsåskådningen. --- Vi stå här vid en punkt som blifvit starkt framhållen af danska filosofer, ända från Holberg till Sibbern och Kierkegaard. Se härom min bok »Danska filosofer« (1909).} Det är en ljusglimt i tidens mörker, att denna sida af personlighetsprincipen mer och mer erkännes. Erkännes den fullständigt, skall en alldeles ny uppfattning af de sista problemen göra sig gällande. Tillsvidare %
% --- p. 70 ---
\opage{70}
kämpa de gamla dogmerna med de nya tviflen. Igenom denna skärseld skall den afgörande segern vinnas för personlighetsprincipen.

\end{document}
